% GENERATED FILE. Edit canonical lessons, then run npm run generate:books.
% canonical-source-hash: fnv1a64:13482fdd55e91232
% canonical-lessons: PA-C205-fine-return, PA-C205-sabh, PA-C205-sabh-review, PA-C205-everything-fine, PA-C205-payoff

\chapter{Everything Is Fine}
\label{ch:pa-a2-sabh}
\hlchaptermodality{205}

\begin{chapteropening}
\textbf{By the end of this chapter:} \emph{I can recognise and use one word for all in a bounded familiar frame.}

Recognise and produce one word for all in a bounded familiar frame through listening, speaking, reading, and writing.
\end{chapteropening}

\section[\texorpdfstring{\pa{ਠੀਕ} \pa{ਹੈ} (thik hai)}{ਠੀਕ ਹੈ (thik hai)}]{Return to a familiar fine-is frame}

\label{lesson:PA-C205-fine-return}



\begin{quote}
\emph{Hear:} \emph{thik hai}; YOU SAY: \textbf{\pa{ਠੀਕ} \pa{ਹੈ}}
\end{quote}

\begin{grammarlens}[title={Grammar Lens}]
Keep these two familiar words ready in this order. The next lesson will own one new word; this opener does not reveal it.
\end{grammarlens}

\subsection*{Your turn}
\begin{itemize}
\raggedright
  \item \emph{Hear:} \emph{fine, is}; YOU SAY: \textbf{\pa{ਠੀਕ} \pa{ਹੈ}}
  \item Say the familiar two-word frame once, then stop.
\end{itemize}

\begin{tcolorbox}[breakable,colback=teal!4,colframe=teal!35!black,title={Before you move on}]
Produce the familiar frame from meaning. This opener owns no new word.
\end{tcolorbox}
\section[\texorpdfstring{\pa{ਸਭ} (sabh)}{ਸਭ (sabh)}]{One whole word for all}

\label{lesson:PA-C205-sabh}



\begin{quote}
Recall the English idea “all,” then listen for one new Punjabi word.
\end{quote}

\subsection*{What to know first}
\begin{quote}
\textbf{\pa{ਸਭ}} — \emph{sabh} — all; everyone; everything here
\end{quote}

Hear and say the whole word before reading it. Keep \textbf{\pa{ਸਭ}} whole. In this chapter it carries one beginner idea: all, understood as everyone or everything in the small frame that follows. The lesson does not open every grammatical use.

\begin{cousinweb}
Punjabi University Patiala's lexical statistics lists \textbf{\pa{ਸਭ}} as an indefinite pronoun with the meanings “all / every / whole.” Its pronoun grammar also lists the word among plural indefinite pronouns and glosses it “all.” Its number introduction independently shows that the same form can serve an indefinite numeral role. Sources: \href{https://www.learnpunjabi.org/statistics.html}{Punjabi University lexical statistics}, \href{https://www.learnpunjabi.org/Pronoun.html}{Punjabi University pronoun grammar}, and \href{https://www.learnpunjabi.org/numbers\_intro.html}{Punjabi University number introduction}.

Those sources support the present meaning, but they do not establish a deeper history. The Punjabi origin remains unknown in the sources used here. Make no historical derivation claim; remember the whole form.
\end{cousinweb}

\subsection*{Your turn}
\begin{itemize}
\raggedright
  \item \emph{Hear:} \emph{sabh}; YOU SAY: \textbf{\pa{ਸਭ}}
  \item \emph{Read it:} \textbf{\pa{ਸਭ}}
  \item \emph{Write it:} copy \textbf{\pa{ਸਭ}} once, checking the whole word
\end{itemize}

\begin{tcolorbox}[breakable,colback=teal!4,colframe=teal!35!black,title={Before you move on}]
Produce the one word for “all.”
\end{tcolorbox}
\section[\texorpdfstring{\pa{ਸਭ} (sabh)}{ਸਭ (sabh)}]{Retrieve the new all word on its own}

\label{lesson:PA-C205-sabh-review}



\begin{quote}
\emph{Hear:} \emph{sabh}; YOU SAY: \textbf{\pa{ਸਭ}}
\end{quote}

\subsection*{Your turn}
\begin{itemize}
\raggedright
  \item \emph{Hear:} \emph{all, everyone, or everything here}; YOU SAY: \textbf{\pa{ਸਭ}}
  \item \emph{Read it:} \textbf{\pa{ਸਭ}}
  \item \emph{Write it:} cover and write \textbf{\pa{ਸਭ}}
\end{itemize}

\begin{grammarlens}[title={Grammar Lens}]
Keep one whole form and one bounded idea here. The familiar two-word frame returns only after this separate retrieval.
\end{grammarlens}

\begin{tcolorbox}[breakable,colback=teal!4,colframe=teal!35!black,title={Before you move on}]
Say and then write the word once from memory.
\end{tcolorbox}
\section[\texorpdfstring{\pa{ਸਭ} \pa{ਠੀਕ} \pa{ਹੈ} (sabh thik hai)}{ਸਭ ਠੀਕ ਹੈ (sabh thik hai)}]{Put all before one familiar frame}

\label{lesson:PA-C205-everything-fine}



\begin{quote}
Retrieve the new word for “all,” then the familiar words for “fine, is.”
\end{quote}

\begin{grammarlens}[title={Grammar Lens}]
\begin{quote}
\textbf{\pa{ਸਭ} \pa{ਠੀਕ} \pa{ਹੈ}} — \emph{sabh thik hai} — everything is fine
\end{quote}

For this one frame, put the whole word \textbf{\pa{ਸਭ}} before familiar \textbf{\pa{ਠੀਕ} \pa{ਹੈ}}. Here “all” is naturally understood as “everything.” Nothing else is added.
\end{grammarlens}

\subsection*{Your turn}
\begin{itemize}
\raggedright
  \item \emph{Hear:} \emph{sabh thik hai}; YOU SAY: \textbf{\pa{ਸਭ} \pa{ਠੀਕ} \pa{ਹੈ}}
  \item \emph{Read it:} \textbf{\pa{ਸਭ} \pa{ਠੀਕ} \pa{ਹੈ}}
  \item \emph{Write it:} cover and write \textbf{\pa{ਸਭ} \pa{ਠੀਕ} \pa{ਹੈ}}
\end{itemize}

\begin{tcolorbox}[breakable,colback=teal!4,colframe=teal!35!black,title={Before you move on}]
Produce the complete “everything is fine” frame.
\end{tcolorbox}
\section[\texorpdfstring{\pa{ਸਭ} \pa{ਠੀਕ} \pa{ਹੈ} (sabh thik hai)}{ਸਭ ਠੀਕ ਹੈ (sabh thik hai)}]{A bounded everything-is-fine payoff}

\label{lesson:PA-C205-payoff}



\begin{quote}
Retrieve the new word for “all” alone before rebuilding its familiar frame.
\end{quote}

\subsection*{Your turn: four skills}
\begin{itemize}
\raggedright
  \item \emph{Hear:} \emph{sabh thik hai}; YOU SAY: \textbf{\pa{ਸਭ} \pa{ਠੀਕ} \pa{ਹੈ}}
  \item \emph{Read it:} \textbf{\pa{ਸਭ} \pa{ਠੀਕ} \pa{ਹੈ}}
  \item \emph{Write it:} cover the line and write \textbf{\pa{ਸਭ} \pa{ਠੀਕ} \pa{ਹੈ}}
\end{itemize}

\begin{tcolorbox}[breakable,colback=teal!4,colframe=teal!35!black,title={Before you move on}]
Say the new word, then produce the complete familiar frame.

This is the twenty-eighth prerequisite tranche exposed by Punjabi A2 mock drafting. It does not unblock Punjabi A2 mock 1 by itself. It is not an A2 mock, not A2 readiness or attainment, and not external or human validation.
\end{tcolorbox}
